
\section{Advanced Activities}

The following activities are intended for strong master's or PhD students.
They emphasize rigorous mathematical formulation of control co-design
problems, including trajectory variables, objective-function transformations,
path and boundary constraints, free-final-time problems, multi-scenario
formulations, and multiphase systems.

\subsection*{Problem 1: Complete Continuous-Time CCD Formulation with Free Final Time}

Consider the controlled mass--spring--damper system
\begin{equation}
m\ddot{q}(t)+c\dot{q}(t)+kq(t)=u(t)+d(t),
\end{equation}
where
\begin{equation}
m=1.5~\mathrm{kg},
\qquad
d(t)=2e^{-0.5t}\sin(3t).
\end{equation}

The plant design variables are
\begin{equation}
0.5\leq k\leq 8,
\qquad
0.1\leq c\leq 4,
\qquad
0.5\leq F_{\max}\leq 6,
\end{equation}
and the final time is also a design variable:
\begin{equation}
2\leq t_f\leq 8.
\end{equation}

The control trajectory satisfies
\begin{equation}
|u(t)|\leq F_{\max}.
\end{equation}

The initial and terminal conditions are
\begin{equation}
q(0)=1,
\qquad
\dot{q}(0)=0,
\end{equation}
\begin{equation}
q(t_f)=0,
\qquad
\dot{q}(t_f)=0.
\end{equation}

The displacement and velocity must satisfy
\begin{equation}
|q(t)|\leq 1.2,
\qquad
|\dot{q}(t)|\leq 2.5.
\end{equation}

Minimize
\begin{equation}
J
=
\int_0^{t_f}
\left[
q(t)^2
+
0.1\dot{q}(t)^2
+
0.02u(t)^2
\right]dt
+
0.03k^2
+
0.02c^2
+
0.04F_{\max}^2
+
0.2t_f.
\end{equation}

\begin{enumerate}
    \item Define a first-order state vector and derive the state equations
    in the form
    \begin{equation}
    \dot{\mathbf{x}}
    =
    \mathbf{f}
    \left(
    \mathbf{x},
    u,
    \mathbf{x}_p,
    t
    \right).
    \end{equation}

    \item Define the complete plant-design vector, control trajectory, state
    trajectory, and free-time variable.

    \item Write the problem in the general Bolza form
    \begin{align}
    \min_{\mathbf{x}_p,u(\cdot),\mathbf{x}(\cdot),t_f}
    \quad
    &
    \Phi
    \left(
    \mathbf{x}(0),
    \mathbf{x}(t_f),
    \mathbf{x}_p,
    t_f
    \right)
    +
    \int_0^{t_f}
    L
    \left(
    \mathbf{x},
    u,
    \mathbf{x}_p,
    t
    \right)dt,
    \\
    \text{subject to}
    \quad
    &
    \dot{\mathbf{x}}
    -
    \mathbf{f}
    \left(
    \mathbf{x},
    u,
    \mathbf{x}_p,
    t
    \right)
    =
    \mathbf{0},
    \end{align}
    together with all path, boundary, and bound constraints.

    \item Classify every constraint as one of the following:
    \begin{enumerate}
        \item dynamic equality constraint,
        \item path inequality constraint,
        \item boundary equality constraint,
        \item simple bound.
    \end{enumerate}

    \item Determine whether $F_{\max}$ is redundant when $u(t)$ is also
    bounded. Explain why it is still a meaningful plant or actuator-design
    variable in this formulation.

    \item Replace the terminal equalities with the terminal tolerance
    \begin{equation}
    \left\|
    \mathbf{x}(t_f)
    \right\|_2
    \leq 10^{-2},
    \end{equation}
    and write the equivalent scalar inequality constraint.

    \item Explain which terms in the objective discourage the optimizer from
    selecting the largest possible actuator, the stiffest possible spring,
    or the longest possible maneuver time.
\end{enumerate}


\subsection*{Problem 2: Conversion among Lagrange, Mayer, and Bolza Forms}

Consider the dynamic optimization problem
\begin{equation}
\min_{u(\cdot)}
\quad
J
=
\int_0^{t_f}
\left[
x(t)^2
+
\rho u(t)^2
\right]dt,
\end{equation}
subject to
\begin{equation}
\dot{x}(t)=a x(t)+b u(t),
\qquad
x(0)=x_0.
\end{equation}

\begin{enumerate}
    \item Introduce an auxiliary state $z(t)$ and convert the Lagrange
    objective into an equivalent Mayer objective.

    \item State the initial condition for $z(t)$ and write the augmented
    dynamics.

    \item Prove that
    \begin{equation}
    z(t_f)
    =
    \int_0^{t_f}
    \left[
    x(t)^2+\rho u(t)^2
    \right]dt.
    \end{equation}

    \item Now consider the Bolza objective
    \begin{equation}
    J
    =
    \frac{1}{2}q_f x(t_f)^2
    +
    \int_0^{t_f}
    \left[
    x(t)^2+\rho u(t)^2
    \right]dt.
    \end{equation}
    Convert it into a pure Mayer problem.

    \item Add the integral energy constraint
    \begin{equation}
    \int_0^{t_f}u(t)^2dt
    \leq E_{\max}.
    \end{equation}
    Introduce another auxiliary state and rewrite the integral inequality
    as a terminal boundary constraint.

    \item Show that the resulting augmented system contains no integral
    objective or integral constraint.

    \item Explain why the Mayer transformation is mathematically exact but
    may still affect numerical conditioning and scaling.

    \item Generalize the transformation to the vector running cost
    \begin{equation}
    \mathbf{L}
    \left(
    \mathbf{x},
    \mathbf{u},
    \mathbf{x}_p,
    t
    \right)
    \in\mathbb{R}^{n_L}.
    \end{equation}
\end{enumerate}


\subsection*{Problem 3: Time Normalization for a Free-Final-Time CCD Problem}

Consider the general free-final-time CCD problem
\begin{equation}
\min_{\mathbf{x}_p,\mathbf{x}_c,\mathbf{u}(\cdot),t_f}
\quad
\Phi
\left(
\mathbf{x}(t_f),
\mathbf{x}_p,
\mathbf{x}_c,
t_f
\right)
+
\int_{t_0}^{t_f}
L
\left(
\mathbf{x},
\mathbf{u},
\mathbf{x}_p,
\mathbf{x}_c,
t
\right)dt,
\end{equation}
subject to
\begin{equation}
\dot{\mathbf{x}}
=
\mathbf{f}
\left(
\mathbf{x},
\mathbf{u},
\mathbf{x}_p,
\mathbf{x}_c,
t
\right).
\end{equation}

Introduce normalized time
\begin{equation}
\tau
=
\frac{t-t_0}{t_f-t_0},
\qquad
0\leq\tau\leq1.
\end{equation}

\begin{enumerate}
    \item Derive
    \begin{equation}
    \frac{dt}{d\tau}=t_f-t_0.
    \end{equation}

    \item Show that the transformed dynamics are
    \begin{equation}
    \frac{d\mathbf{x}}{d\tau}
    =
    (t_f-t_0)
    \mathbf{f}
    \left(
    \mathbf{x},
    \mathbf{u},
    \mathbf{x}_p,
    \mathbf{x}_c,
    t(\tau)
    \right).
    \end{equation}

    \item Transform the integral objective to the fixed interval
    $\tau\in[0,1]$.

    \item Transform the path constraint
    \begin{equation}
    \mathbf{g}
    \left(
    \mathbf{x}(t),
    \mathbf{u}(t),
    \mathbf{x}_p,
    t
    \right)
    \leq\mathbf{0}.
    \end{equation}

    \item Transform the explicit time-dependent disturbance
    \begin{equation}
    d(t)=A\sin(\omega t+\phi).
    \end{equation}

    \item For the scalar system
    \begin{equation}
    \dot{x}=-px+u,
    \qquad
    x(0)=1,
    \qquad
    x(t_f)=0,
    \end{equation}
    write the complete normalized-time formulation when
    \begin{equation}
    J
    =
    t_f
    +
    \int_0^{t_f}
    \left(
    x^2+0.1u^2
    \right)dt.
    \end{equation}

    \item Explain why $t_f$ appears in both the transformed dynamics and the
    transformed running cost.

    \item State the condition that must be imposed to prevent the optimizer
    from selecting
    \begin{equation}
    t_f=t_0.
    \end{equation}
\end{enumerate}


\subsection*{Problem 4: Formulating Actuator Dynamics, Rate Limits, and Energy Limits}

Consider the plant
\begin{equation}
\dot{x}_1=x_2,
\end{equation}
\begin{equation}
m\dot{x}_2=-kx_1-cx_2+f_a+d(t),
\end{equation}
where $f_a(t)$ is the actual actuator force.

The actuator command is $u_c(t)$, but the actuator has first-order dynamics:
\begin{equation}
\dot{f}_a
=
\frac{1}{\tau_a}
\left(
u_c-f_a
\right).
\end{equation}

The actuator capacity and bandwidth are design variables:
\begin{equation}
0.5\leq F_{\max}\leq5,
\qquad
0.02\leq\tau_a\leq0.3.
\end{equation}

The actuator is subject to
\begin{equation}
|u_c(t)|\leq F_{\max},
\end{equation}
\begin{equation}
|f_a(t)|\leq F_{\max},
\end{equation}
\begin{equation}
|\dot{f}_a(t)|\leq R_{\max},
\end{equation}
and
\begin{equation}
\int_0^T
f_a(t)^2dt
\leq E_{\max}.
\end{equation}

\begin{enumerate}
    \item Construct an augmented state vector that includes the actuator
    dynamics.

    \item Introduce an energy-accumulation state $e(t)$ such that
    \begin{equation}
    \dot{e}=f_a^2,
    \qquad
    e(0)=0.
    \end{equation}

    \item Rewrite the energy limit as a terminal boundary constraint.

    \item Rewrite the actuator-rate constraint entirely in terms of
    $u_c$, $f_a$, and $\tau_a$, without using $\dot{f}_a$ explicitly.

    \item Write the complete augmented dynamics in residual form
    \begin{equation}
    \mathbf{r}
    \left(
    \dot{\mathbf{x}},
    \mathbf{x},
    u_c,
    \mathbf{x}_p,
    t
    \right)
    =
    \mathbf{0}.
    \end{equation}

    \item Identify which of the following are plant variables, control
    trajectories, states, path constraints, and boundary constraints:
    \begin{equation}
    k,\quad c,\quad F_{\max},\quad \tau_a,\quad
    u_c(t),\quad f_a(t),\quad e(t).
    \end{equation}

    \item Explain why imposing only
    \begin{equation}
    |u_c(t)|\leq F_{\max}
    \end{equation}
    is not sufficient to guarantee that the actual actuator force is
    physically feasible.

    \item Add the actuator-design penalty
    \begin{equation}
    J_{\mathrm{act}}
    =
    \alpha_FF_{\max}^2
    +
    \frac{\alpha_\tau}{\tau_a},
    \end{equation}
    and explain the physical tradeoff represented by the second term.
\end{enumerate}


\subsection*{Problem 5: Multi-Scenario Robust CCD Formulation}

A controlled system is described under scenario $s$ by
\begin{equation}
\dot{\mathbf{x}}_s
=
\mathbf{f}
\left(
\mathbf{x}_s,
\mathbf{u}_s,
\mathbf{x}_p,
\boldsymbol{\theta}_s,
t
\right),
\qquad
s=1,\ldots,N_s,
\end{equation}
where $\mathbf{x}_p$ is a common plant design and
$\boldsymbol{\theta}_s$ contains scenario-dependent parameters.

Each scenario has objective
\begin{equation}
J_s
=
\Phi_s
\left(
\mathbf{x}_s(t_f),
\mathbf{x}_p
\right)
+
\int_0^{t_f}
L_s
\left(
\mathbf{x}_s,
\mathbf{u}_s,
\mathbf{x}_p,
t
\right)dt.
\end{equation}

The robust objective is
\begin{equation}
J_R
=
\sum_{s=1}^{N_s}p_sJ_s
+
\lambda
\max_{s}
J_s,
\end{equation}
where
\begin{equation}
p_s\geq0,
\qquad
\sum_{s=1}^{N_s}p_s=1.
\end{equation}

\begin{enumerate}
    \item Write the complete multi-scenario CCD problem when the plant
    design $\mathbf{x}_p$ is common to all scenarios but each scenario has
    its own open-loop control trajectory $\mathbf{u}_s(t)$.

    \item Introduce an epigraph variable $\eta$ and replace
    \begin{equation}
    \max_s J_s
    \end{equation}
    with smooth algebraic inequalities.

    \item Write the resulting objective and all epigraph constraints.

    \item Add scenario-wise path constraints
    \begin{equation}
    \mathbf{g}_s
    \left(
    \mathbf{x}_s,
    \mathbf{u}_s,
    \mathbf{x}_p,
    t
    \right)
    \leq\mathbf{0}.
    \end{equation}

    \item Add the common actuator-capacity variable $F_{\max}$ and impose
    \begin{equation}
    |u_s(t)|\leq F_{\max},
    \qquad
    s=1,\ldots,N_s.
    \end{equation}

    \item Now assume that one causal feedback controller
    \begin{equation}
    \mathbf{u}_s(t)
    =
    \boldsymbol{\pi}
    \left(
    \mathbf{y}_s(t);
    \mathbf{x}_c
    \right)
    \end{equation}
    must be shared by all scenarios. Rewrite the decision-variable set.

    \item Explain the difference between scenario-dependent open-loop
    controls and a common feedback law from the viewpoint of information
    availability and implementability.

    \item If each scenario has $n_x$ states, $n_u$ controls, and $N+1$
    transcription nodes, determine the total number of trajectory decision
    variables for:
    \begin{enumerate}
        \item scenario-dependent open-loop controls;
        \item common parameterized feedback with $n_c$ controller parameters.
    \end{enumerate}
\end{enumerate}


\subsection*{Problem 6: Multiphase Hybrid Control Co-Design Formulation}

A robotic system performs two phases:

\begin{enumerate}
    \item a free-motion phase,
    \item a contact phase.
\end{enumerate}

During phase 1,
\begin{equation}
\dot{\mathbf{x}}^{(1)}
=
\mathbf{f}^{(1)}
\left(
\mathbf{x}^{(1)},
\mathbf{u}^{(1)},
\mathbf{x}_p,
t
\right),
\qquad
t\in[t_0,t_s].
\end{equation}

At contact, the state undergoes the jump
\begin{equation}
\mathbf{x}^{(2)}(t_s^+)
=
\boldsymbol{\Delta}
\left(
\mathbf{x}^{(1)}(t_s^-),
\mathbf{x}_p
\right).
\end{equation}

During phase 2,
\begin{equation}
\dot{\mathbf{x}}^{(2)}
=
\mathbf{f}^{(2)}
\left(
\mathbf{x}^{(2)},
\mathbf{u}^{(2)},
\mathbf{x}_p,
t
\right),
\qquad
t\in[t_s,t_f].
\end{equation}

The switching time $t_s$ and final time $t_f$ are decision variables.

The objective is
\begin{align}
J
=&
\int_{t_0}^{t_s}
L^{(1)}
\left(
\mathbf{x}^{(1)},
\mathbf{u}^{(1)},
\mathbf{x}_p,
t
\right)dt
\nonumber\\
&+
\int_{t_s}^{t_f}
L^{(2)}
\left(
\mathbf{x}^{(2)},
\mathbf{u}^{(2)},
\mathbf{x}_p,
t
\right)dt
+
\Phi
\left(
\mathbf{x}^{(2)}(t_f),
\mathbf{x}_p,
t_f
\right).
\end{align}

\begin{enumerate}
    \item Write the full multiphase CCD decision set, including plant
    variables, both state trajectories, both control trajectories, $t_s$,
    and $t_f$.

    \item Write the dynamic equality constraints for both phases.

    \item Write the phase-linkage constraint associated with the jump map.

    \item Add the event condition
    \begin{equation}
    h_{\mathrm{contact}}
    \left(
    \mathbf{x}^{(1)}(t_s^-),
    \mathbf{x}_p
    \right)
    =
    0.
    \end{equation}

    \item Add the time-ordering constraints
    \begin{equation}
    t_0<t_s<t_f.
    \end{equation}

    \item Suppose the normal contact force in phase 2 is
    \begin{equation}
    F_n
    =
    \psi
    \left(
    \mathbf{x}^{(2)},
    \mathbf{u}^{(2)},
    \mathbf{x}_p
    \right).
    \end{equation}
    Write the unilateral-contact path constraint.

    \item Suppose the impact impulse must satisfy
    \begin{equation}
    I
    \left(
    \mathbf{x}^{(1)}(t_s^-),
    \mathbf{x}_p
    \right)
    \leq I_{\max}.
    \end{equation}
    Classify this as a path or boundary constraint and justify your answer.

    \item Convert both phases to fixed normalized-time intervals
    \begin{equation}
    \tau_1,\tau_2\in[0,1].
    \end{equation}
    Derive the scaled dynamics for each phase.

    \item Explain why using one continuous dynamic model without the jump
    map would produce an incorrect formulation.

    \item Describe how the formulation would change if the contact sequence
    itself were an architecture decision rather than fixed in advance.
\end{enumerate}
